\documentclass{article}
\usepackage{amsmath,amssymb,amsthm}
\newtheorem{theorem}{Theorem}
\begin{document}
\begin{theorem}\label{thm:target}
Let $V$ be a finite set and let $\mathcal E\subseteq\binom V4$ have
$m=|\mathcal E|\le 7$. There is a coloring $c:V\to\{0,1\}$ whose
restriction to every $e\in\mathcal E$ takes both values.
\end{theorem}
\begin{proof}
Color each vertex independently and uniformly with one of two colors.
For each edge $e$, let $B_e$ be the event that $e$ is monochromatic.
Since each edge has four vertices, $\Pr(B_e)=2/2^4=1/8$.
For every coloring, the indicator of $\bigcup_{e\in\mathcal E}B_e$
is at most the sum of the indicators of the $B_e$.
Averaging this pointwise inequality over the finite probability space gives
\[
 \Pr\!\left(\bigcup_{e\in\mathcal E}B_e\right)
 \le \sum_{e\in\mathcal E}\Pr(B_e)=\frac{m}{8}.
\]
Hence
\[
 \Pr\!\left(\bigcap_{e\in\mathcal E}B_e^c\right)
 =1-\Pr\!\left(\bigcup_{e\in\mathcal E}B_e\right)
 \ge 1-\frac{m}{8}\ge\frac18>0.
\]
Any coloring in this event proves the conclusion.
\end{proof}
\end{document}
